\subsection{The original theta complex inside the Hochschild trace realization}\label{endpoint-cc_hochschild--the-original-theta-complex-inside-the-hochschild-trace-realization}

Research continuation, 13 September 2026. Source attribution: the split-support arithmetic programme and its representative maps are the owner's construction. The literature comparison uses Connes and Consani (2023), printed pp.~83--101. This note gives written proofs of the comparison and its boundary qualification. It does not claim a formalization of the adele crossed product, Schwartz analysis, sheaf cohomology, or a uniform arithmetic weight estimate.

\subsubsection{1. Fix the objects before comparing them}\label{endpoint-cc_hochschild--fix-the-objects-before-comparing-them}

Use the original spaces

ENDPOINTSOURCE5C05F74730177742310F7FC7SLOT0END

The arithmetic map and its original quotient are

ENDPOINTSOURCE5C05F74730177742310F7FC7SLOT1END

The factor two, the quotient topology, and the original supported coefficient square remain

ENDPOINTSOURCE5C05F74730177742310F7FC7SLOT2END

Here the supported scalar zero is ENDPOINTSOURCE5C05F74730177742310F7FC7SLOT3END, not the absent scalar ENDPOINTSOURCE5C05F74730177742310F7FC7SLOT4END. Every linear comparison below is a comparison of coefficient fibres; its split lift sends ENDPOINTSOURCE5C05F74730177742310F7FC7SLOT5END to ENDPOINTSOURCE5C05F74730177742310F7FC7SLOT6END, and sends external absence to external absence. In particular an amplitude killed by a comparison remains ENDPOINTSOURCE5C05F74730177742310F7FC7SLOT7END. We do not apply ordinary Hochschild homology directly to the split semiring and assert that additive completion preserves its support masks.

\subsubsection{2. The theta map is the actual Hochschild trace}\label{endpoint-cc_hochschild--the-theta-map-is-the-actual-hochschild-trace}

Let

ENDPOINTSOURCE5C05F74730177742310F7FC7SLOT8END

Connes and Consani, Proposition 3.1, identify the kernel of the inclusion-induced map ENDPOINTSOURCE5C05F74730177742310F7FC7SLOT9END with the span of the differences ENDPOINTSOURCE5C05F74730177742310F7FC7SLOT10END. Proposition 3.3 identifies even Schwartz functions with the specified finite-unit-invariant \textbf{image of ENDPOINTSOURCE5C05F74730177742310F7FC7SLOT11END}. This is not an identification with every sector of ENDPOINTSOURCE5C05F74730177742310F7FC7SLOT12END.

Their orbit-trace calculation, Lemma 3.2(ii), gives

ENDPOINTSOURCE5C05F74730177742310F7FC7SLOT13END

Consequently the following is a literal factorization of the original arithmetic map:

ENDPOINTSOURCE5C05F74730177742310F7FC7SLOT14END

The last arrow in (H6) is understood on the displayed subspace, where its trace has the stated strong-Schwartz range; an unrestricted trace on all crossed-product classes is not asserted.

The book's normalized map is

ENDPOINTSOURCE5C05F74730177742310F7FC7SLOT15END

The map ENDPOINTSOURCE5C05F74730177742310F7FC7SLOT16END is an isometry, since its squared norm is exactly ENDPOINTSOURCE5C05F74730177742310F7FC7SLOT17END. This assertion uses the positive half-line. An even full-line realization has its separately retained factor two.

\subsubsection{3. The source Gaussian is exactly a Laplacian preimage}\label{endpoint-cc_hochschild--the-source-gaussian-is-exactly-a-laplacian-preimage}

Write the book's generator as ENDPOINTSOURCE5C05F74730177742310F7FC7SLOT18END, and distinguish its Laplacian from a representative increment:

ENDPOINTSOURCE5C05F74730177742310F7FC7SLOT19END

For the unchanged Gaussian ENDPOINTSOURCE5C05F74730177742310F7FC7SLOT20END, direct differentiation gives

ENDPOINTSOURCE5C05F74730177742310F7FC7SLOT21END

Thus the original source function has the exact homological expression

ENDPOINTSOURCE5C05F74730177742310F7FC7SLOT22END

For example, Mellin integration by parts gives

ENDPOINTSOURCE5C05F74730177742310F7FC7SLOT23END

initially in a common convergence half-plane and then by continuation. Multiplication by ENDPOINTSOURCE5C05F74730177742310F7FC7SLOT24END gives ENDPOINTSOURCE5C05F74730177742310F7FC7SLOT25END, not ENDPOINTSOURCE5C05F74730177742310F7FC7SLOT26END. No change of seed or leading coefficient is involved.

Lemma 4.1(iii) of the chapter gives a bijection ENDPOINTSOURCE5C05F74730177742310F7FC7SLOT27END. The differential operator is continuous. The two moment constraints define a closed Frechet subspace, so the open mapping theorem also gives a continuous inverse. We obtain a topological cochain isomorphism

ENDPOINTSOURCE5C05F74730177742310F7FC7SLOT28END

It commutes with the polynomial action of ENDPOINTSOURCE5C05F74730177742310F7FC7SLOT29END. Moreover ENDPOINTSOURCE5C05F74730177742310F7FC7SLOT30END, so ENDPOINTSOURCE5C05F74730177742310F7FC7SLOT31END. Hence the same comparison on both ordered legs gives

ENDPOINTSOURCE5C05F74730177742310F7FC7SLOT32END

The degree-zero cycle ENDPOINTSOURCE5C05F74730177742310F7FC7SLOT33END is transported, not removed. The fibrewise lifts of (H12)--(H13) preserve the original labels and commute with synchronization into the joint label.

\subsubsection{4. Keep the closure comparison as an arrow}\label{endpoint-cc_hochschild--keep-the-closure-comparison-as-an-arrow}

The quotient in Proposition 4.2 is, on the specified trace image,

ENDPOINTSOURCE5C05F74730177742310F7FC7SLOT34END

There is always a canonical continuous map

ENDPOINTSOURCE5C05F74730177742310F7FC7SLOT35END

Under the prior global-retraction theorem ENDPOINTSOURCE5C05F74730177742310F7FC7SLOT36END, the image ENDPOINTSOURCE5C05F74730177742310F7FC7SLOT37END is closed: it is the range of the continuous idempotent ENDPOINTSOURCE5C05F74730177742310F7FC7SLOT38END, equivalently the kernel of ENDPOINTSOURCE5C05F74730177742310F7FC7SLOT39END. Then (H15) is a topological isomorphism. This use of the prior analytic theorem is explicit; a finite matrix computation does not prove it.

The chapter's full spectral statement concerns ENDPOINTSOURCE5C05F74730177742310F7FC7SLOT40END: the Laplacian eigenvalues are ENDPOINTSOURCE5C05F74730177742310F7FC7SLOT41END for all nontrivial zeros, with their jets. Its critical-only construction later in the chapter is a different quotient. The following calculation supplies a map between the two kinds of coefficient object without presuming injectivity.

\subsubsection{5. A specified critical-line comparison and its exact finite kernel}\label{endpoint-cc_hochschild--a-specified-critical-line-comparison-and-its-exact-finite-kernel}

Choose the plus-sign Mellin restriction

ENDPOINTSOURCE5C05F74730177742310F7FC7SLOT42END

This is continuous: in logarithmic coordinates, ENDPOINTSOURCE5C05F74730177742310F7FC7SLOT43END and all derivatives have faster-than-polynomial decay, with bounds controlled by finitely many defining seminorms of ENDPOINTSOURCE5C05F74730177742310F7FC7SLOT44END. Fourier transformation therefore lands continuously in Schwartz space. The plus sign is deliberate; the chapter's displayed Fourier transform uses the opposite sign, related by ENDPOINTSOURCE5C05F74730177742310F7FC7SLOT45END.

Put

ENDPOINTSOURCE5C05F74730177742310F7FC7SLOT46END

For ENDPOINTSOURCE5C05F74730177742310F7FC7SLOT47END, Mellin factorization gives

ENDPOINTSOURCE5C05F74730177742310F7FC7SLOT48END

The remaining Mellin factor is Schwartz: at the origin ENDPOINTSOURCE5C05F74730177742310F7FC7SLOT49END, and at infinity it is Schwartz; repeated logarithmic differentiation and integration by parts control every Fourier seminorm. The zeta factor is a Schwartz multiplier, as used in the chapter's Section 5. Thus (H18) induces an actual continuous linear map

ENDPOINTSOURCE5C05F74730177742310F7FC7SLOT50END

It also factors through (H15), by continuity and the closed target relation space. Its operator equations are

ENDPOINTSOURCE5C05F74730177742310F7FC7SLOT51END

Theorem 5.4 identifies the target coefficient quotient with global sections of the Scaling Site quotient sheaf. Section 7 below records a zero-mode qualification needed before using the printed raw-section formula at the boundary ENDPOINTSOURCE5C05F74730177742310F7FC7SLOT52END.

Now let ENDPOINTSOURCE5C05F74730177742310F7FC7SLOT53END be a finite packet of actual zeros, with the retained multiplicities, and use the previously constructed inclusion

ENDPOINTSOURCE5C05F74730177742310F7FC7SLOT54END

Chinese remainders give the specified decomposition ENDPOINTSOURCE5C05F74730177742310F7FC7SLOT55END, according to whether ENDPOINTSOURCE5C05F74730177742310F7FC7SLOT56END. Then

ENDPOINTSOURCE5C05F74730177742310F7FC7SLOT57END

\textbf{Proof on the off-critical blocks.} For ENDPOINTSOURCE5C05F74730177742310F7FC7SLOT58END, the function ENDPOINTSOURCE5C05F74730177742310F7FC7SLOT59END is smooth on the real line and all its derivatives are polynomially bounded. It is a continuous Schwartz multiplier, commutes with ENDPOINTSOURCE5C05F74730177742310F7FC7SLOT60END, and preserves both the ideal and its closure. Therefore ENDPOINTSOURCE5C05F74730177742310F7FC7SLOT61END is invertible on ENDPOINTSOURCE5C05F74730177742310F7FC7SLOT62END. If ENDPOINTSOURCE5C05F74730177742310F7FC7SLOT63END, (H20) gives

ENDPOINTSOURCE5C05F74730177742310F7FC7SLOT64END

and hence ENDPOINTSOURCE5C05F74730177742310F7FC7SLOT65END. This kills the full generalized block under this particular map, not just its reduced eigenline.

\textbf{Proof on a critical block.} Let ENDPOINTSOURCE5C05F74730177742310F7FC7SLOT66END have order ENDPOINTSOURCE5C05F74730177742310F7FC7SLOT67END. The continuous derivative evaluations descend to

ENDPOINTSOURCE5C05F74730177742310F7FC7SLOT68END

Indeed these functionals vanish on ENDPOINTSOURCE5C05F74730177742310F7FC7SLOT69END, and hence on its closure. The chain rule gives

ENDPOINTSOURCE5C05F74730177742310F7FC7SLOT70END

The right side is exactly the original Mellin-jet observation. Its composite with ENDPOINTSOURCE5C05F74730177742310F7FC7SLOT71END is the Chinese-remainder projection to the ENDPOINTSOURCE5C05F74730177742310F7FC7SLOT72END-block, by ENDPOINTSOURCE5C05F74730177742310F7FC7SLOT73END. Taking all critical blocks provides a left inverse on ENDPOINTSOURCE5C05F74730177742310F7FC7SLOT74END, proving (H22).

In split coordinates, (H22) says

ENDPOINTSOURCE5C05F74730177742310F7FC7SLOT75END

not external absence. The critical quotient's spectral location does not establish injectivity of (H19). Formula (H22) calculates the exact finite failure of injectivity that a comparison argument would have to address.

\subsubsection{6. The original metric supplies the Laplacian defect, not an assumed Hodge star}\label{endpoint-cc_hochschild--the-original-metric-supplies-the-laplacian-defect-not-an-assumed-hodge-star}

Fix an original representative ENDPOINTSOURCE5C05F74730177742310F7FC7SLOT76END with

ENDPOINTSOURCE5C05F74730177742310F7FC7SLOT77END

The normalized representative ENDPOINTSOURCE5C05F74730177742310F7FC7SLOT78END has the same Gram matrix. With ENDPOINTSOURCE5C05F74730177742310F7FC7SLOT79END acting in multiplicative coordinates,

ENDPOINTSOURCE5C05F74730177742310F7FC7SLOT80END

Thus the shift and the source boundary are both accounted for. In particular applying ENDPOINTSOURCE5C05F74730177742310F7FC7SLOT81END to the entire relation curve ENDPOINTSOURCE5C05F74730177742310F7FC7SLOT82END leaves its Gram matrix, restriction return, determinant ratio, and curvature cost unchanged.

For the tensor-weight algebraic formula, let ENDPOINTSOURCE5C05F74730177742310F7FC7SLOT83END, keep the original ENDPOINTSOURCE5C05F74730177742310F7FC7SLOT84END, and put

ENDPOINTSOURCE5C05F74730177742310F7FC7SLOT85END

Expansion proves

ENDPOINTSOURCE5C05F74730177742310F7FC7SLOT86END

Writing ENDPOINTSOURCE5C05F74730177742310F7FC7SLOT87END gives the equivalent centered energy identity

ENDPOINTSOURCE5C05F74730177742310F7FC7SLOT88END

For (H29), expand ENDPOINTSOURCE5C05F74730177742310F7FC7SLOT89END; the middle ENDPOINTSOURCE5C05F74730177742310F7FC7SLOT90END terms cancel and the remaining expression is ENDPOINTSOURCE5C05F74730177742310F7FC7SLOT91END. For (H30), substitute ENDPOINTSOURCE5C05F74730177742310F7FC7SLOT92END, or expand with ENDPOINTSOURCE5C05F74730177742310F7FC7SLOT93END. These are exact identities before any estimate. The real part of the first term on the right of (H30) is the retained error in the expected negative energy inequality.

The existing rank-two source formula for ENDPOINTSOURCE5C05F74730177742310F7FC7SLOT94END consequently gives a rank-at-most-four formula for ENDPOINTSOURCE5C05F74730177742310F7FC7SLOT95END. This rank fact alone bounds neither its relative norm nor its effect on the growing arithmetic packet.

An actual constituent inclusion ENDPOINTSOURCE5C05F74730177742310F7FC7SLOT96END with ENDPOINTSOURCE5C05F74730177742310F7FC7SLOT97END satisfies ENDPOINTSOURCE5C05F74730177742310F7FC7SLOT98END. No commutation of ENDPOINTSOURCE5C05F74730177742310F7FC7SLOT99END with a metric restriction-loss operator is required.

The full jet formula is

ENDPOINTSOURCE5C05F74730177742310F7FC7SLOT100END

When ENDPOINTSOURCE5C05F74730177742310F7FC7SLOT101END, the nilpotent factor ENDPOINTSOURCE5C05F74730177742310F7FC7SLOT102END is invertible and commutes with ENDPOINTSOURCE5C05F74730177742310F7FC7SLOT103END. Therefore the powers of the nilpotent part have exactly the same kernels as ENDPOINTSOURCE5C05F74730177742310F7FC7SLOT104END: Jordan lengths remain unchanged. In particular negative real diagonal spectrum is not the same as self-adjointness on a positive Hilbert metric. Requiring exact self-adjointness on a multiple-zero packet would also force semisimplicity. The chapter itself retains nontrivial Jordan actions in Remark 5.5.

\subsubsection{7. The boundary zero mode needs to be retained separately}\label{endpoint-cc_hochschild--the-boundary-zero-mode-needs-to-be-retained-separately}

There is a concrete qualification in the printed Section 5 construction. Proposition 5.2(iii), p.~94, places the raw family ENDPOINTSOURCE5C05F74730177742310F7FC7SLOT105END in the smooth-section sheaf vanishing at ENDPOINTSOURCE5C05F74730177742310F7FC7SLOT106END. But the same proof, p.~95, explicitly gives its zeroth normalized Fourier coefficient as

ENDPOINTSOURCE5C05F74730177742310F7FC7SLOT107END

For the original seed (H9) this is

ENDPOINTSOURCE5C05F74730177742310F7FC7SLOT108END

The nonvanishing is exact, not inferred from numerics. In (H11) the gamma factor is nonzero at ENDPOINTSOURCE5C05F74730177742310F7FC7SLOT109END; and ENDPOINTSOURCE5C05F74730177742310F7FC7SLOT110END for real ENDPOINTSOURCE5C05F74730177742310F7FC7SLOT111END, since the paired alternating series for ENDPOINTSOURCE5C05F74730177742310F7FC7SLOT112END is strictly positive. Thus the raw coefficient ENDPOINTSOURCE5C05F74730177742310F7FC7SLOT113END does not extend continuously to zero, let alone vanish to all orders. This calculation concerns the literal raw-section assertion, not the validity of the preceding Hochschild trace and all-zero quotient statements.

A support-retaining repair is to use the augmented Fourier-coordinate object

ENDPOINTSOURCE5C05F74730177742310F7FC7SLOT114END

where flatness is imposed at zero when it belongs to ENDPOINTSOURCE5C05F74730177742310F7FC7SLOT115END. Map a source function to

ENDPOINTSOURCE5C05F74730177742310F7FC7SLOT116END

For ENDPOINTSOURCE5C05F74730177742310F7FC7SLOT117END, reconstruct the full raw Fourier vector by

ENDPOINTSOURCE5C05F74730177742310F7FC7SLOT118END

In the circle's function coordinates, the added constant function has value ENDPOINTSOURCE5C05F74730177742310F7FC7SLOT119END. The different powers in these two coordinate systems are forced by the unitary normalization of the Fourier basis.

The covering transfer satisfies

ENDPOINTSOURCE5C05F74730177742310F7FC7SLOT120END

Accordingly the renormalized zero-line coordinate transforms by ENDPOINTSOURCE5C05F74730177742310F7FC7SLOT121END. The two components in (H34) therefore have genuine compatible ENDPOINTSOURCE5C05F74730177742310F7FC7SLOT122END-actions. Normalized scaling is trivial on the zero line. Nothing has been silently discarded.

\paragraph{The extra zero line is an actual closed boundary}\label{endpoint-cc_hochschild--the-extra-zero-line-is-an-actual-closed-boundary}

Here is an explicit source proof of its fate in the quotient. For ENDPOINTSOURCE5C05F74730177742310F7FC7SLOT123END, form in the original Frechet space

ENDPOINTSOURCE5C05F74730177742310F7FC7SLOT124END

For every fixed Schwartz seminorm the seminorm of the integrand's dilated factor is bounded by a constant times ENDPOINTSOURCE5C05F74730177742310F7FC7SLOT125END. Hence the Gaussian integral converges in Schwartz space. Parity and both original moment conditions persist, so ENDPOINTSOURCE5C05F74730177742310F7FC7SLOT126END. This does not assert uniform source seminorm bounds as ENDPOINTSOURCE5C05F74730177742310F7FC7SLOT127END.

The normalized dilation identity

ENDPOINTSOURCE5C05F74730177742310F7FC7SLOT128END

and Fourier transformation give

ENDPOINTSOURCE5C05F74730177742310F7FC7SLOT129END

On every compact interval ENDPOINTSOURCE5C05F74730177742310F7FC7SLOT130END, all nonzero frequencies have magnitude at least ENDPOINTSOURCE5C05F74730177742310F7FC7SLOT131END. Differentiating their Fourier coefficients produces only finite sums of powers of ENDPOINTSOURCE5C05F74730177742310F7FC7SLOT132END, frequency, and ENDPOINTSOURCE5C05F74730177742310F7FC7SLOT133END, times the Gaussian and derivatives of a Schwartz function. Splitting the Gaussian into two equal factors yields, for each flat-section seminorm of derivative order ENDPOINTSOURCE5C05F74730177742310F7FC7SLOT134END, a bound

ENDPOINTSOURCE5C05F74730177742310F7FC7SLOT135END

To see the uniformity at ENDPOINTSOURCE5C05F74730177742310F7FC7SLOT136END, replace every negative power of ENDPOINTSOURCE5C05F74730177742310F7FC7SLOT137END by a corresponding power of ENDPOINTSOURCE5C05F74730177742310F7FC7SLOT138END, use the second Gaussian to absorb frequency powers, and sum an arbitrarily high negative power of ENDPOINTSOURCE5C05F74730177742310F7FC7SLOT139END. The required constants depend on finitely many Schwartz seminorms of ENDPOINTSOURCE5C05F74730177742310F7FC7SLOT140END.

Thus

ENDPOINTSOURCE5C05F74730177742310F7FC7SLOT141END

Since the last scalar is nonzero, the closed source-generated module in (H34) contains the entire renormalized zero line. Multiplication by local smooth functions gives the corresponding sheaf statement. Once this line is retained, its removal from the quotient is justified by a displayed family of original source relations. Algebraically, if ENDPOINTSOURCE5C05F74730177742310F7FC7SLOT142END is the closed source module and ENDPOINTSOURCE5C05F74730177742310F7FC7SLOT143END, then ENDPOINTSOURCE5C05F74730177742310F7FC7SLOT144END, and the quotient is the nonzero-mode quotient. The split lift makes the zero line a supported boundary, not absence.

This is a written repair and comparison, not a completed Lean proof of the heat-limit or sheaf assertion. It supplies the missing zero-mode convention without adding an artificial third vanishing condition to the original ENDPOINTSOURCE5C05F74730177742310F7FC7SLOT145END.

A separate local typo on p.~100 is also relevant when reading the quotient argument: ENDPOINTSOURCE5C05F74730177742310F7FC7SLOT146END tends to ENDPOINTSOURCE5C05F74730177742310F7FC7SLOT147END as ENDPOINTSOURCE5C05F74730177742310F7FC7SLOT148END, not as ENDPOINTSOURCE5C05F74730177742310F7FC7SLOT149END. We use the former limit. The published page and the supplied scan were inspected; this is not an OCR substitution.

\subsubsection{8. What has and has not been identified}\label{endpoint-cc_hochschild--what-has-and-has-not-been-identified}

The established chain is

ENDPOINTSOURCE5C05F74730177742310F7FC7SLOT150END

where the first comparison is the actual two-term-complex isomorphism (H12), not a degree-forgetting isomorphism of every displayed space. On finite arithmetic packets the last arrow has the calculated kernel (H22), and every surviving critical jet has the explicit recovery (H23).

The polynomial Laplacian uses the actual arithmetic action ENDPOINTSOURCE5C05F74730177742310F7FC7SLOT151END. The curvature restriction-loss operator uses the actual source metrics. They are not identified with one another, and their eigenspaces are not presumed to agree. Equations (H27)--(H31) give the actual connection through the original boundary form ENDPOINTSOURCE5C05F74730177742310F7FC7SLOT152END.

The chapter's Hodge dictionary is explicitly heuristic on p.~91. Its Corollary 4.3 makes a negative-spectrum statement equivalent to RH; it does not prove that statement by declaring the displayed star to be the adjoint for the original Gram metric. Here (H29) computes that adjoint defect instead.

\subsubsection{9. Formal and computational scope}\label{endpoint-cc_hochschild--formal-and-computational-scope}

\texttt{SplitZeroLaplacianControl.lean} formalizes the matrix identities (H29), constituent intertwining, the full quadratic jet (H31), and the general lemma that an intertwiner sends a nilpotent block to zero when the corresponding shifted target operator is injective. The actual target multiplier's invertibility and the critical-jet recovery are the analytic written proof in Section 5, not assumed custom Lean axioms.

The four preceding new modules formalize the relation Gram curve, polynomial local specialization sequence, least invariant hull, and curvature/restriction scalar identities. Their scope and the cyclic gcd/metric-pole consequences are recorded in \texttt{RESEARCH\_NOTE.md}. The strict workflow checks individual sources at \texttt{-\/-trust=0}, treats warnings as errors, and audits the selected 29+9 new declarations against \texttt{propext}, \texttt{Classical.choice}, and \texttt{Quot.sound}. Actual run status is recorded separately after execution; source text alone is not a successful certificate.

The exact Python checks include the original Gaussian differentiation, all factors in (H11), the non-Euclidean polynomial-source Gram metric, generalized jets, the derivative phase ENDPOINTSOURCE5C05F74730177742310F7FC7SLOT153END, and the renormalized covering coefficient. The finite critical/off-critical polynomial calibration is expressly not a set of asserted zeta zeros. Neither those checks nor the Lean algebra prove a uniform arithmetic nonconcentration bound.

\subsubsection{References and reading scope}\label{endpoint-cc_hochschild--references-and-reading-scope}

Connes, A., \& Consani, C. (2023). Hochschild homology, trace map and ENDPOINTSOURCE5C05F74730177742310F7FC7SLOT154END-cycles. In A. Connes, C. Consani, B. I. Dundas, M. Khalkhali, \& H. Moscovici (Eds.), \emph{Cyclic cohomology at 40: Achievements and future prospects} (Proceedings of Symposia in Pure Mathematics, Vol. 105, pp.~83--101). American Mathematical Society. https://doi.org/10.1090/pspum/105/01896

The complete nineteen-page chapter was read in the uploaded volume. The exact trace factor, Laplacian lemma, Hodge dictionary, and zero-mode passage were also inspected as rendered pages. The p.~100 limit and Jordan remark were checked in the author's published offprint. This does not claim reading the full proceedings volume or formal verification of the chapter.

Prior programme inputs: the owner's \texttt{Tau\_Chain\_Descent}, \texttt{Tau\_Global\_Comparison}, \texttt{Tau\_Homotopy\_Boundary\_Control}, and the 13 September 2026 specialization/curvature note; the existing reader's restriction formalization at PR26 commit \texttt{e382e4f61aa4a61b621a8f064dfefc1ab80640d8}. The present branch also retains main commit \texttt{1f7e7c02343884a17df9873566e81a9083c70951} without changing the frozen publication sources.
